% !TEX TS-program = pdflatex
% !TEX encoding = UTF-8 Unicode

% This file is a template using the "beamer" package to create slides for a talk or presentation
% - Giving a talk on some subject.
% - The talk is between 15min and 45min long.
% - Style is ornate.

% MODIFIED by Jonathan Kew, 2008-07-06
% The header comments and encoding in this file were modified for inclusion with TeXworks.
% The content is otherwise unchanged from the original distributed with the beamer package.

\documentclass{beamer}
\setbeamercovered{invisible}



% Copyright 2004 by Till Tantau <tantau@users.sourceforge.net>.
%
% In principle, this file can be redistributed and/or modified under
% the terms of the GNU Public License, version 2.
%
% However, this file is supposed to be a template to be modified
% for your own needs. For this reason, if you use this file as a
% template and not specifically distribute it as part of a another
% package/program, I grant the extra permission to freely copy and
% modify this file as you see fit and even to delete this copyright
% notice. 


\mode<presentation>
{
  \usetheme{Warsaw}
  % or ...

  % or whatever (possibly just delete it)
}


\usepackage[english]{babel}
% or whatever

\usepackage[utf8]{inputenc}
% or whatever

\usepackage{times}
\usepackage[T1]{fontenc}
% Or whatever. Note that the encoding and the font should match. If T1
% does not look nice, try deleting the line with the fontenc.

%%% MATH RELATED

\input{macros.tex}

\title{MATH 350-2 Advanced Calculus} 
\subtitle
{} % (optional)

\author[W.R. Casper] % (optional, use only with lots of authors)
{W.R. Casper}
% - Use the \inst{?} command only if the authors have different
%   affiliation.

\institute[California State University Fullerton] % (optional, but mostly needed)
{
  Department of Mathematics\\
  California State University Fullerton}
% - Use the \inst command only if there are several affiliations.
% - Keep it simple, no one is interested in your street address.

\subject{Talks}
% This is only inserted into the PDF information catalog. Can be left
% out. 



% If you have a file called "university-logo-filename.xxx", where xxx
% is a graphic format that can be processed by latex or pdflatex,
% resp., then you can add a logo as follows:

% \pgfdeclareimage[height=0.5cm]{university-logo}{university-logo-filename}
% \logo{\pgfuseimage{university-logo}}



% Delete this, if you do not want the table of contents to pop up at
% the beginning of each subsection:
\AtBeginSubsection[]
{
  \begin{frame}<beamer>{Outline}
    \tableofcontents[currentsection,currentsubsection]
  \end{frame}
}


% If you wish to uncover everything in a step-wise fashion, uncomment
% the following command: 

%\beamerdefaultoverlayspecification{<+->}


\begin{document}

\begin{frame}
  \titlepage
\end{frame}

\begin{frame}{Outline}
  \tableofcontents
  % You might wish to add the option [pausesections]
\end{frame}

% Since this a solution template for a generic talk, very little can
% be said about how it should be structured. However, the talk length
% of between 15min and 45min and the theme suggest that you stick to
% the following rules:  

% - Exactly two or three sections (other than the summary).
% - At *most* three subsections per section.
% - Talk about 30s to 2min per frame. So there should be between about
%   15 and 30 frames, all told.

\section{Real Analysis Lecture 16}

\subsection{Continuity}

\begin{frame}{Continuity}
Let $(S,d_S)$ and $(T,d_T)$ be metric spaces and $f: S\rightarrow T$ a function.
\pause
\begin{defn}
\pause
Then $f$ is continuous at a point $p\in S$ if 
if for all $\epsilon > 0$ there exists $\delta > 0$ such that
$$d_S(x,p) < \delta\ \Rightarrow\ d_T(f(x),f(p)) < \epsilon.$$
\pause
Equivalently, $f(B_S(p;\delta))\subseteq B_T(f(p);\epsilon)$.
\pause
We say $f$ is continuous on $A\subseteq S$ if $f(x)$ is continuous at every point $p\in A$.
\pause
If $f$ is continuous on $S$, we simply say $f$ is continuous.
\end{defn}
\end{frame}

\begin{frame}{Sequence version}
\pause
\begin{thm}
Let $(S,d_S)$ and $(T,d_T)$ be metric spaces and $f: S\rightarrow T$ a function.
\pause
Then $f$ is continuous at a point $p$ in $S$ if for all sequences $\{x_n\}$ of values in $S$ which converge to $p$, we have
\pause
$$\lim_{n\rightarrow\infty} f(x_n) = f(p).$$
\end{thm}
\pause
Intuitively, continuity means
$$\lim_{n\rightarrow\infty} f(x_n) = f\left(\lim_{n\rightarrow\infty} x_n\right).$$
\end{frame}

\begin{frame}{Compositions}
\begin{thm}
Suppose $(S,d_S)$, $(T,d_T)$, and $(U,d_U)$ are metric spaces and $f: S\rightarrow T$ and $g: T\rightarrow U$ are functions.\\
\pause
If $f$ is continuous at $p$ and $g$ is continuous at $f(p)$, then the composition $h = g\circ f: S\rightarrow U$ defined by $h(x) = g(f(x))$ is continuous at $p$.
\end{thm}
\end{frame}

\begin{frame}{Challenge!}
\begin{prob}
Prove the previous theorem!
\end{prob}
\end{frame}

\begin{frame}{Preimages}
\begin{defn}
Suppose $f: S\rightarrow T$ is a function and that $Y\subseteq T$.\\
\pause
The \textbf{inverse image} or \textbf{preimage} of $Y$ is the set
$$f^{-1}(Y) = \{x: x\in S,\ \text{and}\ f(x)\in Y\}.$$
\end{defn}
\end{frame}

\begin{frame}{Challenge!}
\begin{prob}
Suppose that $f: S\rightarrow T$ and $Y\subseteq Z\subseteq T$.\\
Prove that
$$f^{-1}(Y)\subseteq f^{-1}(Z).$$
\end{prob}
\end{frame}


\begin{frame}{Open sets}
\begin{thm}
Suppose $(S,d_S)$ and $(T,d_T)$ are metric spaces and $f: S\rightarrow T$ is a function.\\
\pause
Then $f$ is continuous if and only if $f^{-1}(V)$ is open for all open subsets $V\subseteq T$.
\end{thm}
\end{frame}

\begin{frame}{Open sets}
\begin{proof}
\pause
$(\Rightarrow)$\\
\pause
Suppose $V\subseteq T$ is open.\\
\pause
Then for all $x\in f^{-1}(V)$, we have $x\in S$ and $f(x)\in V$.\\
\pause
Since $V$ is open, there exists $\epsilon > 0$ such that $B_T(f(x);\epsilon)\subseteq V$.\\
\pause
Since $f$ is continuous at $x$, there exists $\delta > 0$ such that $f(B_S(x;\delta))\subseteq B_T(f(x);\epsilon)\subseteq V$.\\
\pause
Consequently $B_S(x;\delta)\subseteq f^{-1}(V)$.\\
\pause
This shows that $x$ is an interior point of $f^{-1}(V)$.\\
\pause
Since $x$ was arbitrary, this proves $f^{-1}(V)$ is open.
\end{proof}
\end{frame}


\begin{frame}{Open sets}
\begin{proof}
\pause
$(\Leftarrow)$\\
\pause
Suppose that for all open subsets $V\subseteq T$, the preimage $f^{-1}(V)$ is open.\\
\pause
Let $x\in S$.\\
\pause
Then for all $\epsilon > 0$, the ball $B_T(f(x);\epsilon)$ is open.\\
\pause
Therefore the preimage $f^{-1}(B_T(f(x);\epsilon))$ is open.\\
\pause
It contains the point $x$, which must be an interior point.\\
\pause
This means that there exists $\delta > 0$ such that $B_S(x;\delta)\subseteq f^{-1}(B_T(f(x);\epsilon)).$\\
\pause
Thus $f(B_S(x;\delta))\subseteq B_T(f(x);\epsilon)$.
\end{proof}
\end{frame}

\begin{frame}{Closed sets}
\begin{thm}
Suppose $(S,d_S)$ and $(T,d_T)$ are metric spaces and $f: S\rightarrow T$ is a function.\\
\pause
Then $f$ is continuous if and only if $f^{-1}(A)$ is closed for all closed subsets $A\subseteq T$.
\end{thm}
\end{frame}

\begin{frame}{Relatively open sets}
\begin{thm}
Let $(S,d_S)$ and $(T,d_T)$ be metric spaces and $f: S\rightarrow T$ a function.\\
If $f$ is continuous on a subset $A\subseteq S$, then for any open subset $V\subseteq T$ there exists an open subset $U\subseteq S$ with 
$$f^{-1}(V)\cap A = U\cap A.$$
\end{thm}
\end{frame}

\begin{frame}{Relatively open sets}
\begin{proof}
\pause
Consider the subspace $(A,d_A=d_S)$ of $S$.\\
\pause
The function $f$ restricts to a continuous function $f|_A: A\rightarrow T$.\\
\pause
Therefore if $V\subseteq T$ is open, the preimage
$$f|_A^{-1}(V) = \{x: x\in A,\ \text{and}\ f(x)\in V\} = f^{-1}(V)\cap A$$
\pause
is open in the subspace $A$ of $S$.\\
\pause
This is the same thing as
$$ f^{-1}(V)\cap A = U\cap A$$
for some open subset $$U\subseteq S$$.
\end{proof}
\end{frame}

\begin{frame}{Compact sets}
\begin{thm}
Let $(S,d_S)$ and $(T,d_T)$ be metric spaces and $f: S\rightarrow T$ a function.\\
If $f$ is continuous on a compact subset $A\subseteq S$, then $f(A)$ is compact.
\end{thm}
\pause
\begin{proof}
\pause
Suppose that $\{V_i: i\in I\}$ is an open cover of $f(A)$.\\
\pause
Then $f^{-1}(V_i)\cap A = U_i\cap A$ for some open subset $U_i\subseteq S$.\\
\pause
The open subsets $\{U_i: i\in I\}$ cover $A$, and since $A$ is compact there is a finite subcover $\{U_1,\dots,U_n\}$.\\
\pause
This means $A\subseteq \bigcup_{i=1}^n U_i$ and therefore
\pause
$$f(A)\subseteq \bigcup_{i=1}^n V_i.$$
\pause
This proves that $\{V_1,\dots, V_n\}$ is a finite subcover of $f(A)$.
\end{proof}
\end{frame}

\begin{frame}{Extreme value theorem}
\begin{thm}
Suppose $(S,d_S)$ is a metric space and $f: S\rightarrow \mathbb{R}$ is a function (to Euclidean space).\\
If $f$ is continuous on a compact subset $A\subseteq S$, then there exists $a,b\in A$ such that
$$f(a) = \inf f(A),\quad\text{and}\quad f(b) = \sup f(A).$$
\end{thm}
\end{frame}

\begin{frame}{Extreme value theorem}
\pause
\begin{proof}
\pause
The set $f(A)$ is compact in $\mathbb{R}$, so it is closed and bounded.\\
\pause
The Completeness Axiom therefore tells us that it has an infimum $y=\inf(f(A))$ and supremum $z=\sup(f(A))$.\\
\pause
The infimum and supremum of a set are adherent points of that set.\\
\pause
Therefore, since $f(A)$ is closed $y,z\in f(A)$.\\
\pause
This means that there exists $a,b\in A$ such that $y=f(a)$ and $z=f(b)$.
\end{proof}
\end{frame}

\subsection{Connectedness}

\begin{frame}{Characterization of intervals}
\begin{thm}
If $I\subseteq \mathbb{R}$ is an interval, then there does not exist two nonempty, disjoint open sets $U,V\subseteq\mathbb{R}$ (Euclidean space), with
$$I\subseteq U\cup V.$$
\end{thm}
\end{frame}

\begin{frame}{Connected spaces}
\begin{defn}
A metric space $(S,d)$ is called \textbf{disconnected} if there exists two nonempty, disjoint open subsets $U,V\subseteq S$ with $S=U\cup V$.\\
\pause
We call $S$ \textbf{connected} if it is not disconnected.\\
\pause
A subset $A\subseteq S$ is connected if the corresponding subspace $(A,d)$ is connected.
\end{defn}
\pause
Examples:
\begin{itemize}
\pause
\item the real line with the Euclidean metric $(\mathbb{R},d_{eucl})$ is connected
\pause
\item the real line with the discrete metric $(\mathbb{R},d_{disc})$ is connected
\pause
\item euclidean space $\mathbb{R}^n$ is connected
\pause
\item an interval $I\subseteq \mathbb{R}$ with the Euclidean metric is connected
\end{itemize}
\end{frame}

\begin{frame}{Intermediate value theorem}
\begin{thm}[Intermediate Value Theorem]
Suppose $f: [a,b]\rightarrow\mathbb{R}$ is a continuous function with the Euclidean metric.\\
Then for any value $y$ between $f(a)$ and $f(b)$, there exists $c\in [a,b]$ with $f(c) =y$.
\end{thm}
\end{frame}


\end{document}


